\section{Building and verifying the formalization}\label{app:verify}

The Lean project is the directory \leanfile{lean/} next to this paper. It uses Lean and Mathlib, both at version v4.34.1
(\leanfile{lean-toolchain}, \leanfile{lakefile.toml}, \leanfile{lake-manifest.json}), and requires the libraries of the eighth
paper~\cite{EighthPaper2026} and, for the corollaries, of the eleventh paper~\cite{MergeDepth2026} through relative paths; those
libraries require the earlier libraries of the series in the same way.
Keep the repository's directory layout, so that these sibling directories stay in place.

\paragraph{Build.} With \texttt{elan} installed, run in \leanfile{lean/}:
\begin{center}
\texttt{lake exe cache get \&\& lake build}
\end{center}
The first command downloads the compiled Mathlib; the second compiles the library and the parts of the earlier libraries that it
uses. Before the first build, fetch the second paper's pinned external dependency as its README describes
(\texttt{python3 verification/fetch\_openai\_metric.py}, run in that paper's directory).

\paragraph{Axioms.} The file \leanfile{PaperAxioms.lean}, generated from the \LaTeX{} sources of this paper by
\leanfile{tools/lean_index.py}, runs \texttt{\#print axioms} on every declaration of this library named here (\cref{app:index}):
\begin{center}
\texttt{lake env lean PaperAxioms.lean}
\end{center}
Every declaration depends at most on \texttt{propext}, \texttt{Classical.choice} and \texttt{Quot.sound}. No Lean file of this
library uses \texttt{sorry}, \texttt{admit} or \texttt{native\_decide}, declares an \texttt{axiom}, or raises
\texttt{maxHeartbeats}. The script regenerates both the index of \cref{app:index} and \leanfile{PaperAxioms.lean}; with the
option \texttt{--index-only} it writes only the index.

\paragraph{Where things are.} The library \leanfile{lean/BoundedStackDepth/} has six layers, each a group of modules.
\begin{itemize}[leftmargin=1.5em]\raggedright
\item \emph{Statements and foundations:} \leanfile{Statement.lean} (the main statement and its corollary for stack DAGs),
\leanfile{Assembly.lean} (the probabilistic method for label cuts), \leanfile{Basic.lean}, \leanfile{Grid.lean},
\leanfile{BandSums.lean}, \leanfile{Choice.lean}, \leanfile{Walks.lean}, \leanfile{Prominence.lean} and
\leanfile{TopProminence.lean}.
\item \emph{One walk} (\cref{sec:ladders}): \leanfile{Decomposition.lean}, \leanfile{Excursions.lean}, \leanfile{FramesU.lean},
\leanfile{FrameSystem.lean}, \leanfile{FrameDynamics.lean}, \leanfile{FrameWeight.lean}, \leanfile{FrameWeightU.lean},
\leanfile{LadderBands.lean}, \leanfile{CostModel.lean}, \leanfile{HopData.lean}, \leanfile{Events.lean},
\leanfile{PerExcursion.lean} and \leanfile{TheoremA.lean}.
\item \emph{The label} (\cref{sec:label}): \leanfile{Label.lean} and \leanfile{Global.lean}.
\item \emph{Displays} (\cref{sec:displays}): \leanfile{Display.lean}, \leanfile{Revealing.lean}, \leanfile{DisplayTail.lean},
\leanfile{DisplayView.lean}, \leanfile{Entries.lean}, \leanfile{Composition.lean}, \leanfile{ViewInvariance.lean} and
\leanfile{Exposure.lean}.
\item \emph{Accounting and the proof} (\cref{sec:accounting,sec:proof}): \leanfile{Budgets.lean}, \leanfile{EventWeight.lean},
\leanfile{Accounting.lean}, \leanfile{Truncation.lean}, \leanfile{Wrap.lean}, \leanfile{WrapGlobal.lean},
\leanfile{CutDescents.lean}, \leanfile{CutHops.lean}, \leanfile{Totals.lean} and \leanfile{Main.lean}.
\item \emph{Consequences} (\cref{sec:consequences}): \leanfile{DeepNesting.lean}, the first statement of Corollary~N;
\leanfile{Machines.lean}, Corollaries~O and~HV.
\end{itemize}
The root module \leanfile{BoundedStackDepth.lean} imports all of them.
